\documentclass{beamer}
\usepackage{graphicx}
\title{Euro Week 5 Presentation}
\author{Nick, Sriram, Rohit, Silas}
\begin{document}
    \maketitle
    \begin{frame}
        \frametitle{Decision Trees Converging to $E(Y|X)$}

        \begin{itemize}
            \item Think about an n-dim space of all square integrable random variables as vectors. If $|\Omega| = 3$ it's like $R^3$. RVs X have coordinates corresponding to each outcome. 

            \item $E(Y) = \arg \min_c E[(Y - c)^2]$Think about this as $E(Y) = P_{c \vec{1}} Y$ onto the subset of degenerate random variables (constants)
        
            \item $g(X) = E(Y|X)$ is the arg min among all RVs with finite 2nd derivatives that are m.b. function of X of $E[(X-g(X))^2]$ 
            
            \item $E[\{Y - g(X)\}h(X)] = 0, \forall h,$ h measurable bounded function

            \item Assume h can be arbitrarily closely approximated "dense" by sequence $\{s_n\}$, $s_n = \sum_{i=1}^n \alpha_i \mathbbm{1}_{A_i}$ where $A_i$ are events in $\sigma(X)$
        \end{itemize}

        Decision Trees
        \begin{itemize}
            \item Notice that trees are additive models that are essentially simple functions $f(x) = \sum_{m=1}^M c_m \mathbbm{1} (x \in Q_m) $ where M is number of leaf nodes and $c_m $ is the predicted response 

            \item $\{X \in Q_m\}$ is just an event in $\sigma(X)$

            \item Trees directly approximate $E(Y|X)$ via empirical simple functions
        \end{itemize}
    \end{frame}


    \begin{frame}
        \frametitle{XGBoost Poisson Based Objective}

         \small
        \setlength{\abovedisplayskip}{4pt}
        \setlength{\belowdisplayskip}{4pt}
        \setlength{\abovedisplayshortskip}{2pt}
        \setlength{\belowdisplayshortskip}{2pt}
        \begin{columns}[T,onlytextwidth]

        % LEFT COLUMN: Statistical model and likelihood
        
        \begin{column}{0.5\textwidth}

            \includegraphics[width = \textwidth]{Screenshot from 2026-11-17 22-12-53.png}
            \includegraphics[width = \textwidth]{Screenshot from 2026-11-17 22-14-12.png}
            \includegraphics[width = \textwidth]{Screenshot from 2026-11-17 21-42-59.png}
        \end{column}
        % LEFT COLUMN: Statistical model and likelihood
        \begin{column}{0.48\textwidth}
            \textbf{1. Statistical Model}

            \[
                Y_i \mid X_i \sim \operatorname{Poisson}(\mu_i)
            \]

            \[
                p(y_i \mid \mu_i)
                = \frac{\mu_i^{y_i}e^{-\mu_i}}{y_i!}
            \]

            \textbf{2. Prediction on Log Scale}

            \[
                f_i = \log(\mu_i)
                \quad \Longrightarrow \quad
                \mu_i = e^{f_i}
            \]

            The negative log-likelihood is

            \[
                \begin{aligned}
                \ell_i(f_i)
                &= -\log p(y_i \mid \mu_i=e^{f_i})\\
                &= e^{f_i}-y_i f_i+\log(y_i!)
                \end{aligned}
            \]

            Dropping the constant term gives our
            objective contribution for obs i:

            \[
                \boxed{\ell_i(f_i)=e^{f_i}-y_i f_i}
            \]
        \end{column}

        \end{columns}
    \end{frame}


\begin{frame}
    \frametitle{XGBoost: Poisson-Based Objective}

     \small
    \setlength{\abovedisplayskip}{4pt}
    \setlength{\belowdisplayskip}{4pt}
    \setlength{\abovedisplayshortskip}{2pt}
    \setlength{\belowdisplayshortskip}{2pt}
    \begin{columns}[T,onlytextwidth]

        

        % RIGHT COLUMN: Gradients and aggregated objective
        \begin{column}{0.48\textwidth}
            \textbf{3. Gradient and Hessian}
            $g_i = e^{f_i}-y_i = \mu_i-y_i$; 
            $h_i = e^{f_i} = \mu_i$

            \vspace{0.3cm}
            \textbf{4. Aggregation Within a Leaf}

            For observations assigned to leaf \(I_j\), define

            \[
                M = \sum_{i\in I_j} \mu_i,
                \qquad
                Y = \sum_{i\in I_j} y_i
            \]

            A leaf update \(d\) shifts each prediction:
            \(f_i \mapsto f_i+d\).

            The resulting objective is

            \[
                \begin{aligned}
                L(d)
                &= \sum_{i\in I_j}
                \left(e^{f_i+d}-y_i(f_i+d)\right)
                \end{aligned}
            \]

            Ignoring terms independent of \(d\):

            \[
                \boxed{
                L(d)=Me^d-Yd
                }
            \]

            \textbf{This is the leaf-level objective
            being minimized.}
        \end{column}

        \begin{column}{0.48\textwidth}
            \textbf{4. Ignore regularization optimize for d}

            \[
                L'(d)=Me^d-Y 
            \]

            \[
                Me^{d^{*}}-Y = 0
            \]
            
            \[d^* = \log(\frac{Y}{M})\]

            \begin{itemize}
                \item If observed counts \(Y\) exceed predicted counts \(M\), increase the predictions.



                \item If observed counts \(Y\) are below predicted counts \(M\), decrease the predictions.
            \end{itemize}
            
        \end{column}

    \end{columns}
\end{frame}


\end{document}
